% !TEX program = pdflatex
% !TEX encoding = UTF-8
%
% Dokumentacja projektu — Badania Operacyjne (Informatyka 2025/2026)
% Temat: Planer tras turystycznych z algorytmami ABC i Bee Algorithm (taniec pszczół)
%
% Kompilacja:
%   pdflatex main.tex
%   pdflatex main.tex   (drugi przebieg dla spisów i odnośników)
%
\documentclass[11pt,a4paper]{article}

% ── Pakiety podstawowe ───────────────────────────────────────────────────────
\usepackage[utf8]{inputenc}
\usepackage[T1]{fontenc}
\usepackage[polish]{babel}
\usepackage{lmodern}
\usepackage{microtype}

% ── Geometria i nagłówki ─────────────────────────────────────────────────────
\usepackage[a4paper,margin=2.3cm]{geometry}
\usepackage{fancyhdr}
\setlength{\headheight}{14pt}
\pagestyle{fancy}
\fancyhf{}
\fancyhead[L]{Badania Operacyjne 2025/2026}
\fancyhead[R]{Planer tras turystycznych}
\fancyfoot[C]{\thepage}
\renewcommand{\headrulewidth}{0.4pt}

% ── Matematyka, algorytmy, kod ───────────────────────────────────────────────
\usepackage{amsmath,amssymb,amsthm}
\usepackage{mathtools}
\usepackage{float}
\usepackage{algorithm}
\usepackage{algpseudocode}
\usepackage{listings}
\usepackage{xcolor}

\definecolor{codebg}{RGB}{248,248,248}
\definecolor{codekw}{RGB}{0,90,160}
\definecolor{codestr}{RGB}{160,80,0}
\definecolor{codecom}{RGB}{90,120,90}

\lstdefinelanguage{PythonExt}{
    language=Python,
    morekeywords={self,True,False,None,from,as,with,yield},
}

% Definicja jezyka JSON (listings nie ma jej w standardzie)
\lstdefinelanguage{json}{
    morestring=[b]",
    morestring=[d]',
    stringstyle=\color{codestr},
    morekeywords={true,false,null},
    sensitive=true,
    comment=[l]{//},
    commentstyle=\color{codecom}\itshape,
    literate=
        *{0}{{{\color{codekw}0}}}{1}
         {1}{{{\color{codekw}1}}}{1}
         {2}{{{\color{codekw}2}}}{1}
         {3}{{{\color{codekw}3}}}{1}
         {4}{{{\color{codekw}4}}}{1}
         {5}{{{\color{codekw}5}}}{1}
         {6}{{{\color{codekw}6}}}{1}
         {7}{{{\color{codekw}7}}}{1}
         {8}{{{\color{codekw}8}}}{1}
         {9}{{{\color{codekw}9}}}{1}
         {\{}{{{\color{codekw}{\{}}}}{1}
         {\}}{{{\color{codekw}{\}}}}}{1}
         {[}{{{\color{codekw}{[}}}}{1}
         {]}{{{\color{codekw}{]}}}}{1},
}

\lstset{
    backgroundcolor=\color{codebg},
    basicstyle=\ttfamily\footnotesize,
    keywordstyle=\color{codekw}\bfseries,
    stringstyle=\color{codestr},
    commentstyle=\color{codecom}\itshape,
    breaklines=true,
    breakatwhitespace=true,
    showstringspaces=false,
    numbers=left,
    numberstyle=\tiny\color{gray},
    frame=single,
    rulecolor=\color{gray!40},
    captionpos=b,
    tabsize=2,
    columns=fullflexible,
    inputencoding=utf8,
    extendedchars=true,
    literate=%
        {ą}{{\k{a}}}1 {ę}{{\k{e}}}1 {ć}{{\'c}}1 {ł}{{\l{}}}1
        {ń}{{\'n}}1 {ó}{{\'o}}1 {ś}{{\'s}}1 {ź}{{\'z}}1 {ż}{{\.z}}1
        {Ą}{{\k{A}}}1 {Ę}{{\k{E}}}1 {Ć}{{\'C}}1 {Ł}{{\L{}}}1
        {Ń}{{\'N}}1 {Ó}{{\'O}}1 {Ś}{{\'S}}1 {Ź}{{\'Z}}1 {Ż}{{\.Z}}1,
}

% ── Grafika, tabele, listy ───────────────────────────────────────────────────
\usepackage{graphicx}
\usepackage{booktabs}
\usepackage{array}
\usepackage{longtable}
\usepackage{enumitem}
\setlist{itemsep=2pt,topsep=4pt}

% ── Hiperłącza i odnośniki ───────────────────────────────────────────────────
\usepackage[hidelinks,colorlinks=true,linkcolor=black,citecolor=black,urlcolor=blue!60!black]{hyperref}
\usepackage{cleveref}

% ── Polskie tłumaczenia środowiska algorithm/algpseudocode ───────────────────
\floatname{algorithm}{Algorytm}
\renewcommand{\algorithmicrequire}{\textbf{Wejście:}}
\renewcommand{\algorithmicensure}{\textbf{Wyjście:}}
\renewcommand{\algorithmicwhile}{\textbf{dopóki}}
\renewcommand{\algorithmicdo}{\textbf{wykonuj}}
\renewcommand{\algorithmicend}{\textbf{koniec}}
\renewcommand{\algorithmicfor}{\textbf{dla}}
\renewcommand{\algorithmicforall}{\textbf{dla każdego}}
\renewcommand{\algorithmicif}{\textbf{jeżeli}}
\renewcommand{\algorithmicthen}{\textbf{wtedy}}
\renewcommand{\algorithmicelse}{\textbf{w przeciwnym razie}}
\renewcommand{\algorithmicreturn}{\textbf{zwróć}}

% ── Skróty notacyjne ─────────────────────────────────────────────────────────
\newcommand{\R}{\mathbb{R}}
\newcommand{\N}{\mathbb{N}}

% ─────────────────────────────────────────────────────────────────────────────
\title{
    \vspace{-1.5cm}
    \textbf{Badania Operacyjne} \\[2pt]
    \large Informatyka 2025/2026 \\[18pt]
    \textbf{Planer tras turystycznych --- optymalizacja trasy} \\[6pt]
    \large z wykorzystaniem algorytmów \emph{Artificial Bee Colony} oraz \emph{Bee Algorithm}
}
\author{
    \textbf{Skład zespołu:} \\[4pt]
    \begin{tabular}{l}
        Bartosz Tokarz \\
        Adrian Michalski \\
        Mikołaj Kaleta \\
    \end{tabular}
}
\date{\today}

% ═════════════════════════════════════════════════════════════════════════════
\begin{document}

\maketitle
\thispagestyle{empty}

\begin{abstract}
\noindent
Projekt dotyczy zadania planowania trasy turystycznej na siatce komórek z~ważonymi
przejściami, ograniczeniami budżetowymi, czasowymi oraz minimalną i~maksymalną liczbą
atrakcji każdego typu. Dla wygenerowania populacji rozwiązań dopuszczalnych
zaproponowano i~zaimplementowano algorytmy \emph{Artificial Bee Colony} (ABC) oraz
\emph{Bee Algorithm} (BA, taniec pszczół), w~których rozwiązanie reprezentowane jest
jako lista \emph{waypointów} (start, atrakcje w~kolejności, koniec) łączonych krokami
A\textsuperscript{*} na siatce 8-kierunkowej z~twardym zakazem powtarzania krawędzi.
\end{abstract}

\vspace{0.4cm}
\hrule
\tableofcontents
\hrule
\vspace{0.6cm}

% ═════════════════════════════════════════════════════════════════════════════
\section{Wstęp}
\label{sec:wstep}

\subsection*{Cel projektu}

Celem projektu jest opracowanie narzędzia wspomagającego planowanie trasy
turystycznej na zadanej mapie. Turysta wyrusza z~ustalonego punktu \emph{start},
porusza się po siatce komórek o~zróżnicowanych kosztach przejścia i~ma za zadanie
zakończyć podróż w~punkcie \emph{end}, odwiedzając po drodze atrakcje w~taki sposób,
aby zmaksymalizować łączną wartość zdobytych atrakcji przy jednoczesnym spełnieniu
ograniczeń:

\begin{itemize}
    \item budżetowego (łączny koszt biletów wstępu do odwiedzonych atrakcji nie może
          przekroczyć założonego budżetu),
    \item czasowego (łączny czas: ruch po siatce plus czas zwiedzania nie może
          przekroczyć dostępnego budżetu czasu),
    \item dotyczącego liczności typów atrakcji (dla każdego typu zdefiniowane są
          minimalna i~maksymalna liczba atrakcji, które trzeba lub można odwiedzić).
\end{itemize}

Dodatkowo, ze względów modelowych, przyjęto twarde wymaganie, aby trasa
\emph{nie używała dwukrotnie tej samej krawędzi} (przejścia między dwoma sąsiednimi
polami) --- co odpowiada intuicji, że turysta nie chce przechodzić tą samą wąską
ścieżką w~obie strony.

\subsection*{Założenia projektu}

\begin{enumerate}
    \item Mapa jest dyskretną siatką $H \times W$ komórek. W~każdej komórce
          $(r,c)$ zdefiniowana jest waga $w_{r,c} \in \N_{+}$ reprezentująca koszt
          ,,bycia'' w~tej komórce (im wyższa waga, tym trudniejszy/dłuższy teren).
    \item Ruch w~siatce odbywa się w~jednym z~ośmiu kierunków
          (4 osiowe + 4 diagonalne). Koszt przejścia z~komórki $(r,c)$ do sąsiada
          wynosi $w_{r,c}$ pomnożone przez $1{,}0$ dla ruchu osiowego oraz $1{,}4$
          dla ruchu diagonalnego (przybliżenie $\sqrt{2}$).
    \item Atrakcje to wyróżnione komórki o~zadanej wartości, koszcie biletu, czasie
          zwiedzania i~typie. Wizyta polega na wejściu na komórkę atrakcji; ponowne
          odwiedzenie tej samej atrakcji nie daje żadnych dodatkowych punktów.
    \item Algorytm operuje w~paradygmacie \emph{generowania populacji rozwiązań
          dopuszczalnych} --- każde rozwiązanie w~populacji końcowej musi spełniać
          wszystkie ograniczenia (budżet, czas, liczności typów, brak powtórzonych
          krawędzi).
    \item Komponentem dostarczającym rozwiązanie początkowe i~jednocześnie
          mechanizmem poszukiwania jest meta-heurystyka \emph{Artificial Bee Colony}
          (ABC).
    \item Aplikacja składa się z~modułów: \texttt{map\_generator} (generator instancji),
          \texttt{parser}/\texttt{load\_json} (obsługa formatu danych),
          \texttt{abc\_algorithm} (rdzeń algorytmu), \texttt{initial\_population}
          (skrypt uruchomieniowy generujący populację) oraz \texttt{visualize\_paths}
          (wizualizacja). Dodatkowo dostępna jest aplikacja webowa
          (\texttt{app.py}, Streamlit) służąca jako interaktywny edytor mapy.
\end{enumerate}

\subsection*{Zakres dokumentu}

Dokument prezentuje formalny opis problemu i~jego model matematyczny
(rozdział~\ref{sec:zagadnienie}), szczegółowo omawia zaprojektowaną metodę
rozwiązania (rozdział~\ref{sec:algorytmy}), opis aplikacji wraz z~instrukcją
uruchomienia (rozdział~\ref{sec:aplikacja}), zarys planowanych eksperymentów
(rozdział~\ref{sec:eksperymenty}), wnioski i~kierunki rozwoju
(rozdział~\ref{sec:podsumowanie}) oraz literaturę i~dodatek z~podziałem prac.

\section{Opis zagadnienia}
\label{sec:zagadnienie}

\subsection{Sformułowanie problemu}
\label{sec:problem}

Środowisko modelowane jest jako siatka, gdzie każde pole reprezentuje
fragment terenu o~określonym koszcie przejścia.

\begin{itemize}
    \item Każde pole posiada wagę określającą trudność terenu.
    \item Koszt przejścia zależy od~kierunku ruchu:
    \begin{itemize}
        \item ruch prosty (poziomy/pionowy): mnożnik $1$,
        \item ruch po~skosie (diagonalny): mnożnik $1{,}4$.
    \end{itemize}
\end{itemize}

Na planszy rozmieszczone są \textbf{atrakcje turystyczne}:
\begin{itemize}
    \item każda atrakcja ma wartość turystyczną,
    \item każda atrakcja ma koszt odwiedzenia (np.\ cena biletu),
    \item każda atrakcja należy do~określonego \emph{typu} (zabytek,
          przyroda, restauracja itp.).
\end{itemize}

\textbf{Podróżnik:}
\begin{itemize}
    \item startuje w~punkcie $s$,
    \item kończy w~punkcie $t$,
    \item posiada ograniczony czas $T$ oraz budżet $B$.
\end{itemize}

\textbf{Dodatkowe wymagania konstrukcyjne:}
\begin{itemize}
    \item każda atrakcja może być odwiedzona maksymalnie raz,
    \item dla każdego typu atrakcji zdefiniowane są limity dolny $l_{\tau}$
          i~górny $u_{\tau}$ na liczbę atrakcji tego typu w~trasie.
\end{itemize}

Problem stanowi rozszerzenie klasycznego zagadnienia znajdowania ścieżki
o~elementy decyzyjne związane z~wyborem punktów pośrednich --- bliski
\emph{Orienteering Problem} z~ograniczeniami typowymi.

\subsection{Model matematyczny}
\label{sec:model}

\subsubsection{Definicja grafu}

Problem reprezentowany jest jako graf:
\[
G = (V, E)
\]
gdzie:
\begin{itemize}
    \item $V$ --- zbiór pól siatki,
    \item $E$ --- możliwe przejścia (krawędzie) między polami sąsiadującymi
          w~siatce 8-kierunkowej (sąsiedztwo Moore'a).
\end{itemize}

\subsubsection{Koszt przejścia}

Koszt przejścia pomiędzy polami:
\[
c_{ij} = w_i \cdot d_{ij}
\]
gdzie:
\begin{itemize}
    \item $w_i$ --- waga pola $i$,
    \item $d_{ij} \in \{1,\ 1{,}4\}$ --- współczynnik kierunku ruchu
          ($1$ dla ruchu osiowego, $1{,}4$ dla diagonalnego).
\end{itemize}

\subsubsection{Atrakcje}

Niech:
\begin{itemize}
    \item $A \subseteq V$ --- zbiór atrakcji,
    \item $v_i$ --- wartość atrakcji $i$,
    \item $p_i$ --- koszt (cena) atrakcji $i$,
    \item $\tau_i \in \mathcal{T}$ --- typ atrakcji $i$.
\end{itemize}

\subsubsection{Typy atrakcji}

Niech $\mathcal{T}$ oznacza zbiór wszystkich typów atrakcji. Każdemu
typowi $\tau \in \mathcal{T}$ przypisane są następujące parametry:
\begin{itemize}
    \item $t_{\tau}$ --- czas potrzebny na zwiedzenie atrakcji typu $\tau$,
    \item $l_{\tau}$ --- minimalna liczba atrakcji typu $\tau$, które
          należy odwiedzić,
    \item $u_{\tau}$ --- maksymalna liczba atrakcji typu $\tau$, które
          można odwiedzić.
\end{itemize}

\subsubsection{Zmienne decyzyjne}

\[
x_{ij} =
\begin{cases}
1 & \text{jeśli przechodzimy z } i \text{ do } j, \\
0 & \text{w przeciwnym razie,}
\end{cases}
\qquad
y_i =
\begin{cases}
1 & \text{jeśli odwiedzamy atrakcję } i, \\
0 & \text{w przeciwnym razie.}
\end{cases}
\]

\subsubsection{Funkcja celu}

Celem jest maksymalizacja wartości odwiedzonych atrakcji:
\begin{equation}
\max \;\; \sum_{i \in A} v_i \cdot y_i .
\label{eq:obj}
\end{equation}

\subsubsection{Ograniczenia}

\textbf{1.\ Start:}
\begin{equation}
\sum_{j} x_{s j} = 1 .
\label{eq:start}
\end{equation}

\textbf{2.\ Koniec:}
\begin{equation}
\sum_{i} x_{i t} = 1 .
\label{eq:end}
\end{equation}

\textbf{3.\ Zachowanie przepływu:}
\begin{equation}
\sum_{i} x_{i k} = \sum_{j} x_{k j}
\quad \forall k \in V \setminus \{s, t\} .
\label{eq:flow}
\end{equation}

\textbf{4.\ Ograniczenie czasu:}
\begin{equation}
\sum_{(i,j) \in E} c_{ij} \cdot x_{ij}
+ \sum_{i \in A} y_i \cdot t_{\tau_i}
\;\le\; T .
\label{eq:time}
\end{equation}

\textbf{5.\ Ograniczenie budżetu:}
\begin{equation}
\sum_{i \in A} p_i \cdot y_i \;\le\; B .
\label{eq:budget}
\end{equation}

\textbf{6.\ Ograniczenia ilościowe na typy atrakcji:}
\begin{equation}
l_{\tau} \;\le\; \sum_{\substack{i \in A \\ \tau_i = \tau}} y_i \;\le\; u_{\tau}
\quad \forall \tau \in \mathcal{T} .
\label{eq:types}
\end{equation}

\textbf{7.\ Powiązanie atrakcji ze~ścieżką.}
Atrakcja może być odwiedzona tylko wtedy, gdy ścieżka przez nią przechodzi:
\begin{equation}
y_i \;\le\; \sum_{j} x_{j i} \quad \forall i \in A .
\label{eq:link}
\end{equation}

\textbf{8.\ Jednokrotność odwiedzin atrakcji:}
\begin{equation}
y_i \le 1 \quad \forall i \in A
\label{eq:attr-unique}
\end{equation}

\textbf{9.\ Binarność:}
\begin{equation}
x_{ij} \in \{0, 1\}, \quad y_i \in \{0, 1\} .
\label{eq:bin}
\end{equation}

Równania \eqref{eq:obj}--\eqref{eq:bin} formalnie definiują problem
optymalizacji jako mieszany model przepływu w~grafie z~elementami doboru
podzbioru (klasy \emph{Orienteering Problem}).

\subsection{Uwagi modelowe i decyzje projektowe}

\begin{itemize}
    \item \textbf{Złożoność.} Problem jest NP-trudny (zawiera w~sobie
          uogólnione \emph{Orienteering Problem}), co uzasadnia
          zastosowanie metod heurystycznych i~meta-heurystycznych
          --- w~tym projekcie algorytmu \emph{Artificial Bee Colony}
          (patrz~rozdział \ref{sec:algorytmy}).
    \item \textbf{Ścieżka jako konstrukcja pomocnicza.} W~implementacji
          fragmenty trasy pomiędzy waypointami (start, kolejne atrakcje,
          koniec) generowane są algorytmem A\textsuperscript{*}
          (najtańsza ścieżka, koszt zgodny z~$c_{ij}$, sąsiedztwo
          8-kierunkowe). Algorytm uwzględnia dynamiczne zakazy: blokowane
          są pola atrakcji niezaplanowanych w~bieżącej liście waypointów,
          Algorytm A* nie nakłada ograniczeń na ponowne
          wykorzystanie tej samej krawędzi.
    \item \textbf{Atrakcje są blokujące.} W~implementacji pole atrakcji,
          której nie planujemy odwiedzić, jest blokowane dla ścieżki ---
          jeśli komórka atrakcji znajdzie się na trasie, atrakcja zostanie
          ,,zaliczona''. Wymusza to spójność: trasa nie biegnie tranzytem
          przez atrakcje, których nie zamierzaliśmy odwiedzić.
    \item \textbf{Kara vs.\ dopuszczalność.} W~zaproponowanym podejściu nie
          stosujemy funkcji kary za naruszenie ograniczeń: algorytm utrzymuje
          jedynie rozwiązania w~pełni dopuszczalne (spełniające
          \eqref{eq:start}--\eqref{eq:bin}), a~rozwiązania niedopuszczalne
          są odrzucane już na etapie konstrukcji.
    \item \textbf{Kryterium porównawcze.} Do porównywania rozwiązań
          w~populacji ABC stosowana jest leksykograficzna krotka
          \[
              \big(\,
                  \textstyle\sum v_i y_i,\;
                  -(\text{czas trasy}),\;
                  -(\text{koszt ruchu}),\;
                  -(\text{koszt atrakcji})
              \,\big),
          \]
          gdzie minus oznacza preferencję niższej wartości. Pozwala to przy
          jednakowej wartości celu różnicować rozwiązania krótsze czasowo
          i~tańsze, co ułatwia generowanie zróżnicowanej populacji.
\end{itemize}

\section{Opis algorytmów}
\label{sec:algorytmy}

W~ramach projektu rozważano i~częściowo zaimplementowano kilka podejść do
generowania trasy. Głównymi meta-heurystykami, które zostały zaimplementowane
i~służą do generowania populacji rozwiązań dopuszczalnych, są \emph{Artificial
Bee Colony} (ABC) oraz \emph{Bee Algorithm} (BA, tzw.\ taniec pszczół). Dodatkowo,
jako prosty baseline porównawczy i~mechanizm testowania interfejsu,
zaimplementowano dwie pomocnicze procedury:
\emph{Random Walk} oraz \emph{Greedy attractions} (zachłanne dążenie do najbliższej
atrakcji).

\subsection{Artificial Bee Colony (ABC)}
\label{sec:abc}

ABC to populacyjna meta-heurystyka inspirowana zachowaniem pszczół. W~klasycznej
formie wyróżnia trzy typy ,,pszczół'':
\begin{itemize}
    \item \textbf{Pszczoły zatrudnione} (\emph{employed bees}) --- przypisane
          do konkretnych rozwiązań (źródeł pokarmu) i~modyfikujące je lokalnie.
    \item \textbf{Obserwatorzy} (\emph{onlookers}) --- wybierają źródło
          z~prawdopodobieństwem rosnącym z~jakością źródła i~również je modyfikują.
    \item \textbf{Zwiadowcy} (\emph{scouts}) --- zastępują źródła, które od~zbyt
          wielu prób (\emph{trial}) nie poprawiły się, nowymi losowymi rozwiązaniami.
\end{itemize}

Schemat iteracji ABC dla naszego problemu przedstawia~\Cref{alg:abc-overview}.

\begin{algorithm}[H]
\caption{Generowanie populacji metodą ABC (ogólny zarys)}
\label{alg:abc-overview}
\begin{algorithmic}[1]
\Require dane wejściowe \texttt{parsed\_data}, rozmiar populacji $N$, liczba iteracji $I$, limit stagnacji $L$, liczba onlookerów $O$
\Ensure populacja $\mathcal{P} = \{s_1, \dots, s_N\}$ rozwiązań dopuszczalnych
\State $\mathcal{P} \gets \emptyset$
\While{$|\mathcal{P}| < N$}
    \State $s \gets$ \Call{LosoweRozwiazanieDopuszczalne}{}
    \If{$s \ne \bot$ \textbf{i} $s$ nie jest duplikatem w~$\mathcal{P}$}
        \State $\mathcal{P} \gets \mathcal{P} \cup \{s\}$
    \EndIf
\EndWhile
\For{$it = 1, \dots, I$}
    \For{$i = 1, \dots, N$} \Comment{Faza employed}
        \State $s' \gets$ \Call{MutacjaWaypointow}{$\mathcal{P}_i$}
        \If{$s' \ne \bot$ \textbf{i} $s' \succ \mathcal{P}_i$}
            \State $\mathcal{P}_i \gets s'$;\ \  $\mathcal{P}_i.\text{trial} \gets 0$
        \Else
            \State $\mathcal{P}_i.\text{trial} \gets \mathcal{P}_i.\text{trial} + 1$
        \EndIf
    \EndFor
    \State $w \gets$ \Call{WagiRankingowe}{$\mathcal{P}$}
    \For{$k = 1, \dots, O$} \Comment{Faza onlooker}
        \State $i \gets$ \Call{WylosujZWagami}{$w$}
        \State $s' \gets$ \Call{MutacjaWaypointow}{$\mathcal{P}_i$}
        \If{$s' \ne \bot$ \textbf{i} $s' \succ \mathcal{P}_i$}
            \State $\mathcal{P}_i \gets s'$;\ \  $\mathcal{P}_i.\text{trial} \gets 0$
        \Else
            \State $\mathcal{P}_i.\text{trial} \gets \mathcal{P}_i.\text{trial} + 1$
        \EndIf
    \EndFor
    \For{$i = 1, \dots, N$} \Comment{Faza scout}
        \If{$\mathcal{P}_i.\text{trial} \ge L$}
            \State $\mathcal{P}_i \gets$ \Call{LosoweRozwiazanieDopuszczalne}{}
            \State $\mathcal{P}_i.\text{trial} \gets 0$
        \EndIf
    \EndFor
\EndFor
\State \Return $\mathcal{P}$
\end{algorithmic}
\end{algorithm}

Relacja porządku $s' \succ s$ to leksykograficzne porównanie krotek:
\[
\big(\,\mathrm{Value}(s'),\ -\mathrm{Time}(s'),\ -\mathrm{Cost}_{\mathrm{ruch}}(s'),\ -\mathrm{Cost}_{\mathrm{atrakcji}}(s')\,\big)
\;>\;
\big(\,\mathrm{Value}(s),\ -\mathrm{Time}(s),\ -\mathrm{Cost}_{\mathrm{ruch}}(s),\ -\mathrm{Cost}_{\mathrm{atrakcji}}(s)\,\big).
\]

\subsubsection{Reprezentacja rozwiązania}
\label{sec:repr}

Rozwiązanie w~populacji jest dwupoziomowe:

\begin{enumerate}
    \item \textbf{Plan logiczny} --- lista \emph{waypointów}
          $\mathbf{W} = \langle s, a_{i_1}, a_{i_2}, \dots, a_{i_k}, e \rangle$,
          gdzie $s,e$ to punkty startu i~końca, a~$a_{i_j}$ to wybrane atrakcje
          w~zaplanowanej kolejności odwiedzenia.
    \item \textbf{Realizacja fizyczna} --- konkretna ścieżka komórek
          $P = \langle p_0, p_1, \dots, p_L \rangle$, zbudowana z~planu
          logicznego przez łączenie kolejnych waypointów najtańszą ścieżką
          A\textsuperscript{*} z~uwzględnieniem zakazu używania już wykorzystanych
          krawędzi i~niezaplanowanych komórek atrakcji.
\end{enumerate}

Każde rozwiązanie w~populacji przechowuje również:
\begin{itemize}
    \item \texttt{path} --- pełną listę komórek (lista par $(r,c)$),
    \item \texttt{visited\_attractions} --- listę odwiedzonych atrakcji,
    \item \texttt{movement\_cost}, \texttt{attraction\_cost}, \texttt{total\_time}, \texttt{total\_value},
    \item \texttt{type\_counts} --- wektor liczności atrakcji każdego typu,
    \item flagi dopuszczalności: \texttt{budget\_ok}, \texttt{time\_ok}, \texttt{type\_ok}, \texttt{feasible},
    \item \texttt{trial} --- licznik nieudanych modyfikacji (na potrzeby zwiadowców),
    \item \texttt{id}, \texttt{waypoints} --- identyfikator i~plan logiczny.
\end{itemize}

\subsubsection{Wygenerowanie populacji początkowej}
\label{sec:init-pop}

Populacja początkowa składa się z~$N$ losowych rozwiązań \textbf{dopuszczalnych}.
Procedura konstrukcji pojedynczego rozwiązania składa się z~dwóch kroków:

\begin{enumerate}
    \item Losowy wybór listy waypointów spełniający ograniczenia liczności
          typów (\Cref{alg:random-waypoints}).
    \item Materializacja ścieżki przez łączenie kolejnych waypointów
          algorytmem A\textsuperscript{*} z~zakazem powtarzania krawędzi
          i~komórek (\Cref{alg:build-path}). Jeżeli ścieżka się nie udaje
          (np.\ nie ma fizycznej drogi spełniającej ograniczenia),
          rozwiązanie jest odrzucane i~próbujemy następne.
\end{enumerate}

\begin{algorithm}[H]
\caption{\textsc{LosoweWaypointy} --- losowanie planu atrakcji}
\label{alg:random-waypoints}
\begin{algorithmic}[1]
\Require dane mapy, punkty $s, e$, parametr $E_{\max}$ (maksymalna liczba dodatkowych atrakcji)
\Ensure lista waypointów $\mathbf{W}$ lub $\bot$ przy braku spełnienia minimów
\State $\mathrm{chosen} \gets \emptyset$
\ForAll{typ $t$}
    \State $n_t \gets m_t -$ liczba atrakcji typu $t$ obecnych w~$\{s, e\}$
    \If{$n_t > 0$}
        \State wylosuj $n_t$ atrakcji typu $t$ spośród dostępnych
        \If{kandydatów jest za mało} \State \Return $\bot$ \EndIf
        \State dołącz wylosowane do $\mathrm{chosen}$
    \EndIf
\EndFor
\State $E \gets \mathrm{rand}\!\left(0,\ \min(E_{\max},\ |\text{wolnych atrakcji}|)\right)$
\For{$k = 1, \dots, E$}
    \State wylosuj atrakcję $a$ z~prawdopodobieństwem ważonym
           $\frac{v_a + 1}{\tau_{t_a} + 1}$
    \If{$a \notin \mathrm{chosen}$ \textbf{i} dorzucenie nie łamie $M_{t_a}$}
        \State $\mathrm{chosen} \gets \mathrm{chosen} \cup \{a\}$
    \EndIf
\EndFor
\State pomieszaj kolejność $\mathrm{chosen}$
\State \Return konkatenacja $\langle s \rangle$, $\mathrm{chosen}$, $\langle e \rangle$
\end{algorithmic}
\end{algorithm}

\begin{algorithm}[H]
\caption{\textsc{BudujSciezkeZWaypointow} --- konstrukcja ścieżki fizycznej}
\label{alg:build-path}
\begin{algorithmic}[1]
\Require mapa, waypointy $\mathbf{W} = \langle w_0, w_1, \dots, w_K \rangle$
\Ensure ścieżka $P$ (lista komórek) lub $\bot$ jeśli nie istnieje
\State $\mathrm{visited} \gets \{w_0\};\quad \mathrm{usedEdges} \gets \emptyset$
\State $P \gets \langle w_0 \rangle$
\For{$i = 0, \dots, K-1$}
    \State $\text{blokady} \gets (\mathcal{A}\text{-pozycje} \setminus \{w_0, w_K, w_i, w_{i+1}\})\ \cup\ (\mathrm{visited} \setminus \{w_i, w_{i+1}\})$
    \State $S \gets$ \Call{AStar}{$w_i, w_{i+1}$, blokady, $\mathrm{usedEdges}$}
    \If{$S = \bot$} \Return $\bot$ \EndIf
    \For{każda kolejna krawędź $(u,v)$ w~$S$}
        \State $\mathrm{usedEdges} \gets \mathrm{usedEdges} \cup \{\{u,v\}\}$
    \EndFor
    \For{każdy kolejny węzeł $p$ w~$S$ poza pierwszym}
        \If{$p \in \mathrm{visited}$} \Return $\bot$ \EndIf
        \State $\mathrm{visited} \gets \mathrm{visited} \cup \{p\}$;\ \  dołącz $p$ do~$P$
    \EndFor
\EndFor
\State \Return $P$
\end{algorithmic}
\end{algorithm}

Procedura A\textsuperscript{*} jest standardowa --- używa heurystyki
optymistycznej $h(p, g) = w_{\min} \cdot (1{,}4\cdot d_{\mathrm{diag}} +
d_{\mathrm{ax}})$, gdzie $w_{\min}$ to minimalna waga na siatce,
$d_{\mathrm{diag}}$ to liczba kroków diagonalnych, a~$d_{\mathrm{ax}}$ to
liczba kroków osiowych w~ruchu Chebysheva. Zapewnia to dopuszczalność
heurystyki i~optymalność znalezionej trasy w~sensie kosztu ruchu (przy
zadanych blokadach komórek i~zakazach krawędzi).

\subsubsection{Operatory modyfikacji rozwiązania (sąsiedztwo)}
\label{sec:operators}

Modyfikacja rozwiązania (\textsc{MutacjaWaypointow}) polega na zmianie
listy waypointów, a~następnie próbie materializacji nowej ścieżki.
Wybór operatora jest losowy z~rozkładem przedstawionym
w~\Cref{tab:operators}. W~każdej próbie sprawdzane są ograniczenia
liczności typów; jeżeli kandydat nie da się zmaterializować lub nie spełnia
ograniczeń, próba kończy się fiaskiem.

\begin{table}[H]
\centering
\caption{Operatory mutacji listy waypointów (sąsiedztwo).}
\label{tab:operators}
\begin{tabular}{p{3.3cm} c p{8.3cm}}
\toprule
\textbf{Operator} & \textbf{P} & \textbf{Opis} \\
\midrule
Dodaj atrakcję      & 0{,}30 & Wybierany kandydat $a \notin \mathbf{W}$
                                ze zbioru 20 najlepiej rokujących wg
                                $(v_a + 1) / (\tau_{t_a} + 1)$; wstawiany
                                w~losowej pozycji, jeśli nie łamie $M_{t_a}$. \\
Usuń atrakcję       & 0{,}25 & Losowo wybierany element $\mathbf{W}$, którego
                                usunięcie nie złamie $m_{t}$. \\
Zamień sąsiadujące  & 0{,}20 & Wybierane dwie pozycje w~$\mathbf{W}$
                                i~zamieniane miejscami. \\
2-opt (odwróć fragment) & 0{,}20 & Losowo wybierany fragment $[i, j]$
                                i~odwracany. \\
Przesuń atrakcję    & 0{,}05 & Wyjęcie pojedynczego elementu i~wstawienie
                                w~innym miejscu (\emph{shift/relocate}). \\
\bottomrule
\end{tabular}
\end{table}

\subsubsection{Selekcja typu onlooker (wagi rankingowe)}
\label{sec:onlookers}

W~fazie onlookerów wybór modyfikowanego źródła odbywa się na podstawie
\emph{rangi} rozwiązania w~populacji (a~nie surowych wartości funkcji celu),
co zapobiega dominacji jednego źródła. Procedura wag rankingowych jest
następująca:

\begin{enumerate}
    \item Posortuj populację malejąco wg krotki porządkowej (patrz
          \Cref{sec:abc}).
    \item Rozwiązaniu o~randze $r$ (od~0) przypisz wagę $w_r = N - r$.
    \item Losuj indeks źródła z~rozkładem proporcjonalnym do $w_r$.
\end{enumerate}

Liczba prób onlookerów w~iteracji wynosi $O$ (parametr; domyślnie
$O = N$).

\subsubsection{Faza zwiadowcza (scout)}
\label{sec:scouts}

Jeśli licznik nieudanych prób ($\mathrm{trial}$) dla rozwiązania przekroczy
limit $L$, rozwiązanie jest zastępowane nowym losowym rozwiązaniem
dopuszczalnym (\Cref{alg:random-waypoints} + \Cref{alg:build-path}). Po
udanym zastąpieniu licznik prób się zeruje. Jeżeli nie udaje się
wygenerować nowego rozwiązania w~zadanej liczbie prób, licznik również się
zeruje, by nie zapętlać próby zastąpienia.

\subsubsection{Pseudokod operatora mutacji}
\label{sec:mutate-pseudo}

\begin{algorithm}[H]
\caption{\textsc{MutacjaWaypointow}}
\label{alg:mutate}
\begin{algorithmic}[1]
\Require rozwiązanie $s$, $\mathrm{maxAttempts}$
\Ensure rozwiązanie $s'$ (jeśli udało się znaleźć dopuszczalnego sąsiada) lub $\bot$
\For{$j = 1, \dots, \mathrm{maxAttempts}$}
    \State $\mathbf{W}' \gets s.\mathbf{W}$
    \State $\rho \gets \mathrm{rand}(0, 1)$
    \If{$\rho < 0{,}30$} \Comment{Dodaj}
        \State dobierz kandydata $a$ ze zbioru top-20 wg $(v_a+1)/(\tau_{t_a}+1)$
        \If{liczność typu $t_a$ po dodaniu $\le M_{t_a}$}
            \State wstaw $a$ w~losowe miejsce $\mathbf{W}'$
        \Else \State \textbf{continue}
        \EndIf
    \ElsIf{$\rho < 0{,}55$} \Comment{Usuń}
        \State $R \gets$ pozycje, których usunięcie nie złamie $m_{t}$
        \If{$R = \emptyset$} \textbf{continue} \EndIf
        \State usuń element o~losowej pozycji z~$R$
    \ElsIf{$\rho < 0{,}75$} \Comment{Zamień}
        \If{$|\mathbf{W}'| < 2$} \textbf{continue} \EndIf
        \State wylosuj $i, k$; zamień $\mathbf{W}'[i] \leftrightarrow \mathbf{W}'[k]$
    \ElsIf{$\rho < 0{,}95$} \Comment{2-opt}
        \If{$|\mathbf{W}'| < 4$} \textbf{continue} \EndIf
        \State wylosuj $i \le k$; odwróć $\mathbf{W}'[i..k]$
    \Else \Comment{Przesuń}
        \If{$|\mathbf{W}'| < 2$} \textbf{continue} \EndIf
        \State wylosuj $i$, $k$; przenieś $\mathbf{W}'[i]$ na pozycję $k$
    \EndIf
    \State $P \gets$ \Call{BudujSciezkeZWaypointow}{$\mathbf{W}'$}
    \If{$P \ne \bot$ \textbf{i} ocena $P$ daje rozwiązanie dopuszczalne}
        \State \Return rozwiązanie zbudowane z~$\mathbf{W}'$ i~$P$
    \EndIf
\EndFor
\State \Return $\bot$
\end{algorithmic}
\end{algorithm}

\subsubsection{Funkcja oceny ścieżki}
\label{sec:eval}

Funkcja \texttt{evaluate\_path} (plik \texttt{initial\_population.py}) liczy:
\begin{itemize}
    \item \texttt{movement\_cost} --- $\sum_{l} w_{p_l}\cdot\delta_l$,
    \item \texttt{attraction\_cost} --- suma $k_i$ po unikalnych odwiedzonych atrakcjach,
    \item \texttt{total\_time} --- $\texttt{movement\_cost} + \sum \tau_{t_i}$,
    \item \texttt{total\_value} --- $\sum v_i$ po unikalnych odwiedzonych atrakcjach,
    \item \texttt{type\_counts} --- wektor $\big[\sum_{i: t_i = t} \mathbb{1}_{a_i \in \mathcal{V}(P)}\big]_{t = 0, \dots, T-1}$,
    \item flagi: \texttt{budget\_ok} ($\le B$), \texttt{time\_ok} ($\le T_{\max}$),
          \texttt{type\_ok} ($m_t \le \cdot \le M_t$) i~\texttt{feasible} (koniunkcja).
\end{itemize}

Spójność z~modelem matematycznym jest zapewniona: zwracana krotka cech
odpowiada wprost ograniczeniom \eqref{eq:budget}--\eqref{eq:types}.

\subsection{Bee Algorithm (BA, taniec pszczół)}
\label{sec:bee}

Bee Algorithm (BA) to populacyjna meta-heurystyka inspirowana sposobem, w~jaki
pszczoły \emph{tańcem} komunikują jakość znalezionych źródeł pokarmu i~rekrutują
inne osobniki do dalszej eksploracji otoczenia. W~naszej implementacji BA działa
na \textbf{tej samej reprezentacji rozwiązania} (waypointy + materializacja
A\textsuperscript{*}) i~na \textbf{tym samym sąsiedztwie} (operatory mutacji) co
ABC, a~różni się przede wszystkim mechanizmem \emph{rekrutacji} oraz
\emph{porzucania} rozwiązań.

Intuicyjnie odpowiada to następującym rolom:
\begin{itemize}
    \item \textbf{Rekruterzy (tańczące pszczoły)} --- podzbiór najlepszych rozwiązań,
          których jakość jest komunikowana i~które przyciągają kolejne próby
          ulepszania.
    \item \textbf{Zbieraczki (foragers)} --- wykonują lokalne przeszukiwanie,
          generując sąsiadów (mutacje waypointów) wokół rozwiązań-rekruterów.
    \item \textbf{Zwiadowcy (scouts)} --- wprowadzają różnorodność, zastępując część
          słabych / zastanych rozwiązań nowymi losowymi rozwiązaniami dopuszczalnymi.
\end{itemize}

Schemat iteracji BA dla naszego problemu przedstawia~\Cref{alg:bee-overview}.

\begin{algorithm}[H]
\caption{Generowanie populacji metodą BA (ogólny zarys)}
\label{alg:bee-overview}
\begin{algorithmic}[1]
\Require dane wejściowe \texttt{parsed\_data}, rozmiar populacji $N$, liczba iteracji $I$, prawdopodobieństwo rekrutacji $p_r$, udział zwiadowców $r_s$, próg jakości tańca $Q$ (\%), cierpliwość porzucenia $A$
\Ensure populacja $\mathcal{P} = \{s_1, \dots, s_N\}$ rozwiązań dopuszczalnych
\State $\mathcal{P} \gets$ $N$ losowych rozwiązań dopuszczalnych (por.\ \Cref{sec:init-pop})
\State $s^\star \gets \bot$
\For{$it = 1, \dots, I$}
    \State $b \gets \max\limits_{s \in \mathcal{P}} s$ \Comment{najlepsze wg relacji $\succ$}
    \If{$s^\star = \bot$ \textbf{lub} $b \succ s^\star$} \State $s^\star \gets b$ \EndIf
    \State $\theta \gets \mathrm{Value}(s^\star) \cdot (Q/100)$ \Comment{próg jakości tańca}
    \State $E \gets \{ i : \mathrm{Value}(\mathcal{P}_i) \ge \theta \}$ \Comment{rekruterzy}
    \If{$E = \emptyset$} \State $E \gets \{\arg\max_i \mathcal{P}_i\}$ \EndIf
    \State $w_i \gets \mathrm{Value}(\mathcal{P}_i) + 1$ dla $i \in E$ \Comment{wagi rekrutacji}
    \State $F \gets \max(1, \lfloor N \cdot p_r \rfloor)$ \Comment{liczba zbieraczek}
    \For{$k = 1, \dots, F$}
        \State $i \gets$ \Call{WylosujZWagami}{$E, w$}
        \State $s' \gets$ \Call{MutacjaWaypointow}{$\mathcal{P}_i$} \Comment{por.\ \Cref{sec:operators}}
        \If{$s' \ne \bot$ \textbf{i} $s' \succ \mathcal{P}_i$}
            \State $\mathcal{P}_i \gets s'$;\ \ $\mathcal{P}_i.\text{trial} \gets 0$
        \Else
            \State $\mathcal{P}_i.\text{trial} \gets \mathcal{P}_i.\text{trial} + 1$
        \EndIf
    \EndFor
    \State $S \gets \max(1, \lfloor N \cdot r_s \rfloor)$ \Comment{liczba zwiadowców}
    \State wybierz $S$ indeksów o~największym \texttt{trial}
    \ForAll{wybrane $i$}
        \State $\mathcal{P}_i \gets$ \Call{LosoweRozwiazanieDopuszczalne}{}
        \State $\mathcal{P}_i.\text{trial} \gets 0$
    \EndFor
    \For{$i = 1, \dots, N$} \Comment{porzucanie stagnujących źródeł}
        \If{$\mathcal{P}_i.\text{trial} \ge A$}
            \State $\mathcal{P}_i \gets$ \Call{LosoweRozwiazanieDopuszczalne}{}
            \State $\mathcal{P}_i.\text{trial} \gets 0$
        \EndIf
    \EndFor
\EndFor
\State \Return $\mathcal{P}$
\end{algorithmic}
\end{algorithm}

Relacja porządku $s' \succ s$ jest taka sama jak w~ABC (leksykograficzne
porównanie krotek jakości; por.\ \Cref{sec:abc}).

\subsubsection{Reprezentacja rozwiązania}

Reprezentacja (waypointy + materializacja ścieżki) jest identyczna jak w~ABC
(\Cref{sec:repr}).

\subsubsection{Wygenerowanie populacji początkowej}

Inicjalizacja populacji polega na~wygenerowaniu $N$ losowych rozwiązań
\textbf{dopuszczalnych} przez losowanie waypointów i~materializację ścieżki
(\Cref{sec:init-pop}, \Cref{alg:random-waypoints}, \Cref{alg:build-path}).

\subsubsection{Operatory modyfikacji (sąsiedztwo)}

Lokalne przeszukiwanie wykonywane przez zbieraczki wykorzystuje ten sam operator
\textsc{MutacjaWaypointow} co ABC (\Cref{sec:operators}, \Cref{tab:operators},
\Cref{alg:mutate}).

\subsubsection{Rekrutacja (taniec) i dobór obszarów eksploatacji}

W~każdej iteracji wyznaczany jest próg jakości tańca $Q$ (w~\%), a~następnie do
zbioru rekruterów trafiają te rozwiązania, których wartość
$\mathrm{Value}(s)$ jest nie mniejsza niż $Q\%$ najlepszej wartości
zaobserwowanej dotychczas. Zbieraczki losują następnie źródła do eksploracji
z~rozkładu proporcjonalnego do $\mathrm{Value}(s)+1$ (żeby uniknąć wag zerowych).
Parametr $p_r$ steruje liczbą prób ulepszania w~iteracji.

\subsubsection{Zwiadowcy i porzucanie rozwiązań}

W~implementacji zastosowano dwa mechanizmy dywersyfikacji:
\begin{enumerate}
    \item \textbf{Zwiadowcy} --- w~każdej iteracji część populacji (udział $r_s$)
          jest zastępowana nowymi losowymi rozwiązaniami dopuszczalnymi; jako
          kandydatów wybieramy te osobniki, które mają największy licznik
          stagnacji \texttt{trial}.
    \item \textbf{Porzucanie} --- jeżeli licznik \texttt{trial} przekroczy
          \texttt{abandon\_patience} $=A$, źródło jest uznawane za nieperspektywiczne
          i~również zastępowane nowym losowym rozwiązaniem.
\end{enumerate}

\subsection{Podejścia pomocnicze}
\label{sec:baselines}

W~kodzie znajdują się także uproszczone metody pomocnicze, używane głównie
do testowania interfejsu wizualizacji w~aplikacji Streamlit.

\subsubsection{Random Walk}

Algorytm losowy spaceruje od~$s$ po sąsiedztwie 8-kierunkowym, dopóki nie
zostanie wykorzystany budżet kosztu lub nie zostanie osiągnięty punkt $e$.
Nie pilnuje ograniczeń typów, nie unika powtarzania komórek/krawędzi
i~służy wyłącznie jako baseline. Implementacja: \texttt{src/path\_solver.py},
funkcja \texttt{solve\_random\_walk}.

\subsubsection{Greedy attractions}

W~każdym kroku wybierana jest najbliższa (w~metryce Chebysheva) jeszcze
nieodwiedzona atrakcja i~algorytm wykonuje kroki w~jej kierunku, dopóki
budżet pozwala. Implementacja: \texttt{src/path\_solver.py}, funkcja
\texttt{solve\_greedy\_attractions}. Także baseline.

\subsection{Inne pomysły rozważane w~fazie projektowej}
\label{sec:other-ideas}

W~trakcie projektu rozważano kilka alternatywnych dróg, które nie weszły
do finalnej implementacji:
\begin{itemize}
    \item \textbf{Algorytm genetyczny (GA)} --- z~populacją chromosomów
          jako permutacji wybranych atrakcji; krzyżowania typu PMX/OX,
          mutacje swap/insert/2-opt. Architektura aplikacji (\texttt{app.py})
          zakłada możliwość podłączenia GA przez stabilny interfejs
          \texttt{solve(map, params) → path}, więc rozszerzenie jest naturalne.
    \item \textbf{Programowanie liniowe całkowitoliczbowe (ILP)} --- model
          \eqref{eq:obj}--\eqref{eq:bin} można sprowadzić do mieszanej
          formulacji ILP z~przepływami w~grafie siatki; rozważano użycie
          solvera (np.\ CBC) do małych instancji jako wartość referencyjna.
          Decyzja: pominięto ze względu na koszt implementacji i~ograniczoną
          skalowalność.
    \item \textbf{Symulowane wyżarzanie (SA)} --- jako pojedyncze-rozwiązaniowa
          alternatywa, działająca na tej samej reprezentacji waypointów
          i~tym samym zbiorze operatorów. Pomysł porzucono, ponieważ ABC
          dawało już zadowalające wyniki, a~równoległa populacja jest
          wartością samą w~sobie (np.\ do dalszego ranking-u/eksploracji).
    \item \textbf{Ant Colony Optimization (ACO)} --- pomysł na klasycznej
          stronie konstrukcji ścieżki, jednak nieoczywista adaptacja do
          twardych ograniczeń typów; pominięto.
\end{itemize}

W~niniejszej dokumentacji szczegółowo opisane są algorytmy ABC oraz Bee Algorithm,
ponieważ oba podejścia zostały zaimplementowane i~wykorzystywane w~aplikacji.

\section{Aplikacja --- dokumentacja użytkownika}
\label{sec:aplikacja}

Aplikacja jest zestawem narzędzi w~języku Python (3.10+), uruchamianych z~wiersza
poleceń, oraz pomocniczej aplikacji webowej zbudowanej w~bibliotece Streamlit.
Dwie ścieżki pracy uzupełniają się:

\begin{itemize}
    \item \textbf{Tryb CLI} --- generacja mapy, generacja populacji metodą ABC,
          walidacja i~wizualizacja.
    \item \textbf{Tryb web (Streamlit)} --- interaktywny edytor/generator mapy,
          podgląd, prostsze algorytmy (greedy/random) i~moduł porównawczy wyników.
\end{itemize}

\subsection{Wymagania i instalacja}

\paragraph{Wymagania:} Python $\ge 3.10$, biblioteki:
\texttt{numpy}, \texttt{matplotlib}, \texttt{pandas}, \texttt{streamlit}.
Pełna lista znajduje się w~pliku \texttt{requirements.txt}:

\begin{lstlisting}[language=bash, caption={Zawartość pliku \texttt{requirements.txt}}]
streamlit>=1.35.0
pandas>=2.0.0
matplotlib>=3.8
numpy>=1.26
scipy>=1.11
\end{lstlisting}

Instalacja w~środowisku wirtualnym:
\begin{lstlisting}[language=bash]
python -m venv .venv
# Windows:
.venv\Scripts\activate
# Linux/Mac:
source .venv/bin/activate

pip install -r requirements.txt
\end{lstlisting}

\subsection{Charakterystyka danych}
\label{sec:dane}

\subsubsection{Format mapy --- CLI (\texttt{map.json})}

Mapa generowana skryptem \texttt{map\_generator.py} zapisywana jest do pliku
JSON. Schemat:

\begin{lstlisting}[language=json, caption={Schemat \texttt{map.json}}]
{
  "weight": [[int, ...], ...],   // siatka H x W wag (int >= 1)
  "time":    int,                // budzet czasu T_max
  "budget":  int,                // budzet pieniezny B
  "attractions": [
    [row, col, value, cost, type],  // atrakcja: 5 intow
    ...
  ],
  "attraction_types": [
    [time, min, max],            // typ: (czas zwiedzania, min, max)
    ...
  ]
}
\end{lstlisting}

\paragraph{Uwaga.} Atrakcje są listami pięciu intów; \texttt{type} to
indeks (0-based) typu w~tablicy \texttt{attraction\_types}.

\subsubsection{Format mapy --- aplikacja web (scenarios/*.json)}

Aplikacja Streamlit posługuje się rozszerzonym formatem (m.in.\ z~nazwą,
nazwami typów, polami startu/końca jako słowniki):

\begin{lstlisting}[language=json, caption={Schemat \texttt{scenarios/*.json} (wycinek z \texttt{small.json})}]
{
  "name": "Small 8x8 - 4 atrakcje",
  "width": 8, "height": 8,
  "weights": [[3,2,4,...], ...],
  "attractions": [
    {"id": "a1", "x": 1, "y": 5, "value": 5, "type": "zabytek"},
    ...
  ],
  "attraction_types": ["zabytek", "przyroda"],
  "start": {"x": 0, "y": 0},
  "end":   {"x": 7, "y": 7},
  "budget": 30,
  "seed": 7
}
\end{lstlisting}

\subsubsection{Format populacji (\texttt{initial\_population.json})}

\begin{lstlisting}[language=json]
[
  {
    "id": 0,
    "path": [[r, c], [r, c], ...],
    "visited_attractions": [[r, c], ...],
    "used_attractions":    [[r, c], ...],
    "movement_cost": float,
    "attraction_cost": int,
    "total_time": float,
    "total_value": int,
    "type_counts": [int, ...],
    "budget_ok": bool, "time_ok": bool, "type_ok": bool,
    "feasible": bool
  },
  ...
]
\end{lstlisting}

\subsection{Uruchamianie --- tryb CLI}

\subsubsection{Krok 1. Wygenerowanie mapy}

Skrypt \texttt{map\_generator.py} po uruchomieniu wprost generuje mapę
$30 \times 30$ z~trzema typami atrakcji i~zapisuje wynik do \texttt{map.json}:

\begin{lstlisting}[language=bash]
python map_generator.py
\end{lstlisting}

Domyślne wartości (do edycji w~bloku \texttt{if \_\_name\_\_ == "\_\_main\_\_"}):
\begin{itemize}
    \item rozmiar: $30 \times 30$,
    \item zakres wag pól: $[7, 12]$,
    \item budżet czasowy: $800$, budżet pieniężny: $300$,
    \item trzy typy atrakcji z~różnymi parametrami $(\tau, m, M)$
          i~zakresami wartości oraz kosztów.
\end{itemize}

\subsubsection{Krok 2. Generowanie populacji metodą ABC}

\begin{lstlisting}[language=bash]
python initial_population.py \
    --population-size 50 \
    --iters 200 \
    --limit 60 \
    --out initial_population.json
\end{lstlisting}

\paragraph{Parametry skryptu:}
\begin{description}[leftmargin=2.6cm,style=nextline]
    \item[\texttt{-{}-map}] ścieżka do pliku mapy JSON (domyślnie \texttt{map.json}).
    \item[\texttt{-{}-out}] plik wyjściowy populacji (domyślnie \texttt{initial\_population.json}).
    \item[\texttt{-{}-population-size}] rozmiar populacji $N$ (domyślnie 50).
    \item[\texttt{-{}-iters}] liczba iteracji ABC (domyślnie 200).
    \item[\texttt{-{}-limit}] limit stagnacji $L$ uruchamiający fazę zwiadowczą (domyślnie 60).
    \item[\texttt{-{}-onlookers}] liczba prób onlookerów na iterację (domyślnie $= N$).
    \item[\texttt{-{}-seed}] ziarno generatora liczb losowych (opcjonalne).
\end{description}

\paragraph{Wyjście:} plik JSON z~populacją oraz wypisane na stdout
statystyki populacji (\cref{lst:stats}) i~,,najlepsza trasa''.

\begin{lstlisting}[caption={Przykład wyjścia statystyk populacji (\texttt{print\_population\_stats}).},label=lst:stats]
=== STATYSTYKI POPULACJI POCZATKOWEJ ===
Liczba tras: 50
Spelnia wszystko: 50/50
Spelnia budzet: 50/50
Spelnia limit czasu: 50/50
Spelnia ograniczenia typow: 50/50

Sredni calkowity czas: ...
Sredni koszt ruchu: ...
Sredni uzyty budzet: ...
Sredni a liczba atrakcji: ...
Srednia wartosc atrakcji: ...
Srednia dlugosc sciezki (liczba pol): ...

Srednia liczba atrakcji kazdego typu:
  Typ 0: ... (min=..., max=..., time=...)
  ...
\end{lstlisting}

\subsubsection{Krok 3. Walidacja wygenerowanej populacji}

Skrypt \texttt{examples/validate\_population.py} sprawdza, czy:
\begin{itemize}
    \item każda trasa ma flagę \texttt{feasible=true},
    \item żadna krawędź (nieskierowana) nie występuje w~trasie dwa razy.
\end{itemize}

\begin{lstlisting}[language=bash]
python examples/validate_population.py initial_population.json
\end{lstlisting}

Zwraca kod wyjścia 0 w~razie pomyślnej walidacji.

\subsubsection{Krok 4. Wizualizacja tras}

\begin{lstlisting}[language=bash]
python visualize_paths.py
\end{lstlisting}

Domyślnie funkcja \texttt{draw\_paths} czyta \texttt{map.json}
i~\texttt{initial\_population.json}. Można wybrać podzbiór tras po id:

\begin{lstlisting}[language=Python]
from visualize_paths import draw_paths
draw_paths(selected_ids=[0, 1, 5, 10])
\end{lstlisting}

Lub wskazać inny plik populacji:

\begin{lstlisting}[language=bash]
python -c "from visualize_paths import draw_paths; draw_paths(population_file='initial_population_abc.json')"
\end{lstlisting}

Wynik: okno Matplotlib z~wizualizacją (atrakcje jako czarne punkty,
ścieżki w~różnych kolorach, start jako zielony kwadrat, koniec jako
czerwony X).

\subsection{Uruchamianie --- aplikacja web (Streamlit)}

\begin{lstlisting}[language=bash]
streamlit run app.py
\end{lstlisting}

Aplikacja udostępnia dwa tryby (wybierane w~bocznym pasku):

\subsubsection{Tryb ,,Edytor / Generator''}

W~lewej kolumnie znajduje się formularz generatora mapy oraz edytor
parametrów (rozmiar, liczba atrakcji, ziarno, zakresy wartości i~wag,
typy atrakcji, rozkład wag --- \texttt{uniform}/\texttt{clusters},
budżet, punkty startu i~końca). Po wygenerowaniu można dalej ręcznie
edytować mapę (nazwę, budżet, listę atrakcji w~tabeli) oraz wczytywać
i~zapisywać scenariusze (\texttt{scenarios/*.json}).

W~prawej kolumnie widoczna jest mapa z~wagami w~kolorach. Można uruchomić
jeden z~dwóch baselineów:
\begin{itemize}
    \item \textbf{Greedy (zachłanny)} --- \texttt{solve\_greedy\_attractions},
    \item \textbf{Random walk} --- \texttt{solve\_random\_walk}.
\end{itemize}

Pod wynikiem prezentowane są metryki: koszt, budżet, wartość, liczba
odwiedzonych atrakcji oraz informacja, czy ścieżka mieści się w~budżecie.

\subsubsection{Tryb ,,Wyniki''}

Pozwala wczytać plik wyników o~strukturze:

\begin{lstlisting}[language=json]
{
  "map_file": "scenarios/small.json",
  "runs": [
    {
      "algo": "GA",
      "path": [[x, y], ...],
      "history": [fitness_0, fitness_1, ...],
      "fitness": liczba
    },
    ...
  ]
}
\end{lstlisting}

W~zakładkach prezentowane są: \emph{mapa z~nałożonymi ścieżkami},
\emph{wykres konwergencji} (na podstawie \texttt{history}) i~\emph{tabela
zbiorcza} (\texttt{algorithm}, fitness, koszt, wartość, liczba atrakcji,
długość ścieżki, flaga \emph{feasible}).

\subsection{Ustawianie parametrów instancji}

Parametry instancji ustawiane są w~zależności od~ścieżki użycia:

\begin{itemize}
    \item \textbf{CLI} --- bezpośrednio w~pliku \texttt{map\_generator.py}
          (zmienne \texttt{rows}, \texttt{cols}, \texttt{time\_limit},
          \texttt{budget}, wywołania \texttt{add\_attraction\_type}),
          a~następnie uruchomienie \texttt{python map\_generator.py}.
    \item \textbf{Web} --- formularz w~lewej kolumnie trybu
          ,,Edytor / Generator'' (rozmiar, liczba atrakcji, zakresy
          wartości i~wag, typy, dystrybucja, start/koniec, budżet, seed).
          Można też wczytać gotowy scenariusz z~\texttt{scenarios/} lub
          z~uploadu pliku.
\end{itemize}

\subsection{Ustawianie parametrów algorytmu (ABC)}

Parametry sterujące zachowaniem ABC są dostępne jako argumenty linii
poleceń (patrz wcześniejszy podrozdział). W~praktyce:
\begin{itemize}
    \item Zwiększenie \texttt{-{}-iters} i~\texttt{-{}-onlookers} daje lepszą
          eksploatację, ale wydłuża czas.
    \item Niższy \texttt{-{}-limit} sprzyja eksploracji (częstsze użycie
          zwiadowców), wyższy --- eksploatacji.
    \item \texttt{-{}-seed} zapewnia powtarzalność wyników.
\end{itemize}

W~kodzie funkcja \texttt{generate\_abc\_population} przyjmuje też dodatkowy
parametr \texttt{deduplicate} (domyślnie \texttt{True}), pilnujący, by
populacja końcowa nie zawierała duplikatów ścieżek.

\subsection{Wyniki i sposób ich interpretacji}

\paragraph{Statystyki agregowane.} Funkcja
\texttt{print\_population\_stats} wypisuje liczność rozwiązań spełniających
poszczególne ograniczenia (budżet, czas, typy, wszystko), średnie wartości
metryk i~rozkład liczności typów. Pomaga to ocenić ,,zdrowie'' populacji
(np.\ niski odsetek \texttt{feasible} może oznaczać zbyt rygorystyczne
ograniczenia instancji).

\paragraph{Najlepsza trasa.} Skrypt wyświetla cechy najlepszego rozwiązania
wg leksykograficznej krotki porządkowej (Value $\uparrow$, Time $\downarrow$,
movement cost $\downarrow$, attraction cost $\downarrow$).

\paragraph{Wizualizacja.} W~oknie \texttt{visualize\_paths.py}:
\begin{itemize}
    \item czarne punkty --- atrakcje (każda atrakcja niezależnie od typu),
    \item linie kolorowe --- poszczególne trasy z~populacji,
    \item zielony kwadrat --- start,
    \item czerwony X --- koniec.
\end{itemize}

Im więcej linii ,,konwerguje'' do podobnego kształtu, tym silniejsza była
selekcja w~ABC (mniejsza różnorodność końcowa). Jeśli linie są mocno
różne i~rozproszone, populacja jest zróżnicowana --- co bywa pożądane jako
materiał wejściowy dla dalszej optymalizacji.

\paragraph{Walidacja.}
Skrypt \texttt{examples/validate\_population.py} sygnalizuje:
\begin{itemize}
    \item \texttt{bad\_no\_repeated\_edges} --- liczba tras łamiących zakaz
          powtarzania krawędzi (oczekiwane: 0),
    \item \texttt{infeasible} --- liczba rozwiązań niedopuszczalnych
          (oczekiwane: 0).
\end{itemize}
Sumarycznie wypisuje status \texttt{OK} lub \texttt{NOT OK}.

\section{Eksperymenty}
\label{sec:eksperymenty}

W~bieżącej wersji projektu wdrożony jest mechanizm zbierania metryk
populacji (\texttt{print\_population\_stats} w~\texttt{initial\_population.py})
oraz walidator \texttt{examples/validate\_population.py}, co stanowi
podstawę pod planowane eksperymenty. Pełna seria pomiarów porównawczych
nie została jeszcze przeprowadzona --- w~tym rozdziale prezentujemy
jedynie zarys planu eksperymentalnego oraz dane, jakie są możliwe do
uzyskania z~aktualnego kodu.

\subsection{Plan eksperymentalny}

\paragraph{Zmienne niezależne.}
\begin{itemize}
    \item Rozmiar mapy: $H \times W \in \{10\!\times\!10,\ 20\!\times\!20,\ 30\!\times\!30\}$.
    \item Gęstość atrakcji: liczba atrakcji $\in \{4, 15, 30\}$.
    \item Rozkład wag: \texttt{uniform} vs.\ \texttt{clusters}.
    \item Parametry ABC: $N \in \{20, 50, 100\}$, $I \in \{50, 200, 500\}$,
          $L \in \{30, 60, 100\}$, $O \in \{N/2, N, 2N\}$.
    \item Ziarno losowości \texttt{seed} --- minimum 5 powtórzeń dla każdej
          kombinacji.
\end{itemize}

\paragraph{Zmienne zależne (mierzone).}
\begin{itemize}
    \item Frakcja rozwiązań dopuszczalnych w~populacji końcowej.
    \item Średnia, mediana i~maksimum \texttt{total\_value} w~populacji.
    \item Średni \texttt{movement\_cost}, \texttt{total\_time}, \texttt{attraction\_cost}.
    \item Średnia długość ścieżki (liczba pól).
    \item Liczba unikalnych ścieżek (różnorodność populacji).
    \item Czas wykonania (sekundy) jako funkcja $N$, $I$, $O$.
\end{itemize}

\paragraph{Hipotezy do weryfikacji.}
\begin{enumerate}
    \item Zwiększenie $I$ (liczby iteracji) podnosi medianę \texttt{total\_value},
          ale z~malejącym przyrostem (krzywa nasycenia).
    \item Faza onlookerów wzmacnia eksploatację --- wyższe $O$ przy stałym $I$
          poprawia średnią wartość celu, ale zmniejsza różnorodność populacji.
    \item Wyższy limit stagnacji $L$ zwiększa eksploatację kosztem
          eksploracji; bardzo niskie $L$ powoduje nadmierne ,,resetowanie''
          przez zwiadowców i~uniemożliwia dopracowanie rozwiązań.
    \item Rozkład wag \texttt{clusters} (klastry niskokosztowe) daje wyższe
          wartości \texttt{total\_value} w~tym samym budżecie czasu niż
          \texttt{uniform}.
\end{enumerate}

\subsection{Porównanie z baselineami}

W~kodzie dostępne są dwa proste algorytmy referencyjne:
\texttt{solve\_random\_walk} i~\texttt{solve\_greedy\_attractions}
(\texttt{src/path\_solver.py}). Porównawcza tabela ma postać:

\begin{table}[H]
\centering
\caption{Szablon zestawienia wyników (do uzupełnienia).}
\label{tab:results-template}
\begin{tabular}{lcccc}
\toprule
Algorytm & \texttt{value} & \texttt{move\_cost} & \texttt{total\_time} & \texttt{feasible} \\
\midrule
Random Walk        & -- & -- & -- & -- \\
Greedy attractions & -- & -- & -- & -- \\
ABC (best in pop)  & -- & -- & -- & -- \\
ABC (mean in pop)  & -- & -- & -- & -- \\
\bottomrule
\end{tabular}
\end{table}

\subsection{Zbieranie wyników}

Aktualny stan kodu pozwala uzyskać następujące dane bez modyfikacji
implementacji:

\begin{itemize}
    \item Surowe statystyki (\Cref{lst:stats} z~\Cref{sec:aplikacja}) ---
          wypisywane na stdout po uruchomieniu \texttt{initial\_population.py}.
    \item Plik \texttt{initial\_population.json} z~wszystkimi rozwiązaniami
          --- pozwala policzyć dowolne dodatkowe statystyki ex post
          (np.\ w~notatniku).
    \item Wizualizację (\texttt{visualize\_paths.py}) dla wybranych \texttt{id}.
\end{itemize}

\subsection{Wyniki wstępne}

Pełne wyniki eksperymentalne zostaną dołączone w~kolejnej wersji
dokumentacji. Bieżąca implementacja przeszła testy poprawności:
\begin{itemize}
    \item Wszystkie wygenerowane rozwiązania w~populacji końcowej spełniają
          wszystkie ograniczenia (\texttt{feasible=true}) --- co jest
          gwarantowane konstrukcyjnie przez algorytm.
    \item Walidator \texttt{validate\_population.py} potwierdza brak
          powtórzeń krawędzi w~każdej trasie.
    \item Wizualizacja prezentuje sensowne kształty tras (krzywe biegnące
          przez atrakcje, omijające zablokowane atrakcje niezaplanowane).
\end{itemize}

\section{Podsumowanie}
\label{sec:podsumowanie}

\subsection{Wnioski}

\begin{enumerate}
    \item Postawiony problem łączy aspekty kombinatoryczne (wybór podzbioru
          atrakcji z~ograniczeniami typów i~budżetów) z~grafowymi
          (istnienie i~koszt fizycznej ścieżki na siatce 8-kierunkowej
          z~zakazem powtórzeń krawędzi). Tradycyjne sformułowania
          \emph{Orienteering Problem} stanowią dobrą inspirację, ale
          twarde wymaganie braku powtórzeń krawędzi czyni nasz model
          wyraźnie trudniejszym.
    \item Dwupoziomowa reprezentacja rozwiązania (waypointy + materializacja
          A\textsuperscript{*}) okazała się praktyczna --- pozwala
          rozdzielić problem ,,co odwiedzić'' od~,,jak dojść''
          i~drastycznie zawęża przestrzeń sąsiedztwa.
    \item Algorytm ABC daje populację rozwiązań w~pełni dopuszczalnych,
          co eliminuje konieczność funkcji kary i~upraszcza analizę.
    \item Leksykograficzna krotka porządkowa
          ($\mathrm{Value}\uparrow$, $\mathrm{Time}\downarrow$,
          $\mathrm{Cost}_{\mathrm{ruch}}\downarrow$,
          $\mathrm{Cost}_{\mathrm{atrakcji}}\downarrow$) jest tanim, ale
          skutecznym kryterium różnicowania rozwiązań i~kierowania
          eksploatacją.
\end{enumerate}

\subsection{Napotkane problemy}

\begin{itemize}
    \item \textbf{Trudność generacji rozwiązań dopuszczalnych.} W~instancjach
          o~ciasnym budżecie czasowym lub bardzo dużej liczbie atrakcji
          częstość uzyskiwania ścieżek dopuszczalnych przy losowych
          waypointach jest mała --- konieczne jest podbicie liczby prób
          (\texttt{max\_tries}) w~\texttt{\_random\_feasible\_solution}.
    \item \textbf{Twardy zakaz powtórzeń krawędzi.} Wprowadzenie tego
          ograniczenia powoduje, że pewne kombinacje waypointów po prostu
          nie da się fizycznie zrealizować (segment domyka się dopiero
          po długim ,,objeździe''); musi to być uwzględnione w~mutacjach
          jako naturalny kanał odrzuceń.
    \item \textbf{Koszt A\textsuperscript{*} przy każdej mutacji.}
          Każda zaakceptowana mutacja oznacza materializację co najmniej
          jednego segmentu ścieżki; przy populacji 50 i~200 iteracjach
          koszt ten kumuluje się szybko --- optymalizacja A\textsuperscript{*}
          (np.\ z~cachowaniem segmentów lub heurystyką niższej jakości,
          ale szybszą) jest naturalnym kierunkiem przyspieszania.
\end{itemize}

\subsection{Kierunki rozwoju}

\begin{enumerate}
    \item \textbf{Pełna integracja ABC z~aplikacją Streamlit.}
          Architektura \texttt{app.py} przewiduje stabilny interfejs
          \texttt{solve(map, params) → path}; dorzucenie wrappera nad
          \texttt{generate\_abc\_population} pozwoli porównywać ABC
          z~baselineami w~UI.
    \item \textbf{Drugi algorytm meta-heurystyczny.} GA / SA / Tabu Search
          na tej samej reprezentacji waypointów (operatory już mamy)
          dałyby porównanie wewnątrz tej samej klasy metod.
    \item \textbf{Solver ILP referencyjny.} Dla małych instancji
          (np.\ $10 \times 10$, 6--8 atrakcji) sformułowanie z~rozdziału
          \ref{sec:model} można zlecić solverowi --- daje ,,prawdziwie
          optymalną'' wartość celu i~pozwala obiektywnie ocenić jakość
          ABC.
    \item \textbf{Bardziej wyrafinowana funkcja oceny.} Wprowadzenie
          metryki Pareto-typu (niezdominowana wartość vs.\ czas) zamiast
          leksykografii, by jawnie zwracać front Pareto.
    \item \textbf{Cache segmentów A\textsuperscript{*}.} Ponieważ ABC
          modyfikuje pojedyncze pozycje waypointów, większość segmentów
          ścieżki pozostaje niezmienna --- ich przeliczanie można
          zaoszczędzić.
    \item \textbf{Większa baza scenariuszy testowych.} Aktualnie dostępne
          są dwa scenariusze webowe (\texttt{small.json},
          \texttt{medium.json}); benchmark wymaga większej i~bardziej
          zróżnicowanej kolekcji.
\end{enumerate}

\section{Literatura}
\label{sec:literatura}

\begin{thebibliography}{99}

\bibitem{karaboga2005}
D.~Karaboga,
\emph{An Idea Based on Honey Bee Swarm for Numerical Optimization},
Technical Report TR06, Erciyes University, 2005.

\bibitem{karaboga2007}
D.~Karaboga, B.~Basturk,
\emph{A powerful and efficient algorithm for numerical function
optimization: Artificial Bee Colony (ABC) algorithm},
Journal of Global Optimization, 39(3), 2007.

\bibitem{vansteenwegen2011}
P.~Vansteenwegen, W.~Souffriau, D.~Van~Oudheusden,
\emph{The Orienteering Problem: A survey},
European Journal of Operational Research, 209(1), 2011.

\bibitem{hart1968}
P.~E.~Hart, N.~J.~Nilsson, B.~Raphael,
\emph{A Formal Basis for the Heuristic Determination of Minimum Cost Paths},
IEEE Transactions on Systems Science and Cybernetics, 4(2), 1968.

\bibitem{python}
Dokumentacja Pythona 3, \url{https://docs.python.org/3/}.

\bibitem{streamlit}
Streamlit Documentation, \url{https://docs.streamlit.io/}.

\bibitem{matplotlib}
J.~D.~Hunter,
\emph{Matplotlib: A 2D Graphics Environment},
Computing in Science \& Engineering, 9(3), 2007;
\url{https://matplotlib.org/}.

\bibitem{numpy}
C.~R.~Harris et~al.,
\emph{Array programming with NumPy},
Nature, 585, 2020;
\url{https://numpy.org/}.

\end{thebibliography}

\section{Dodatek --- podział pracy}
\label{sec:dodatek}

\noindent
Poniższa tabela prezentuje szczegółowy podział zadań w~projekcie wraz
z~procentowym udziałem członków zespołu. Ze względu na charakter
przedmiotu (Badania Operacyjne) podział został tak skonstruowany, aby
każdy z~członków zespołu uczestniczył w~realnej pracy algorytmicznej
w~porównywalnym wymiarze, przy czym \textbf{rdzeń metaheurystyki ABC}
pozostaje głównym obszarem odpowiedzialności jednego autora.

\begin{longtable}{p{4.4cm} p{5.6cm} p{4.4cm}}
\toprule
\textbf{Etap realizacji} & \textbf{Zakres pracy} & \textbf{Osoba / \% udział} \\
\midrule
\endfirsthead

\toprule
\textbf{Etap realizacji} & \textbf{Zakres pracy} & \textbf{Osoba / \% udział} \\
\midrule
\endhead

Model matematyczny problemu
& Formalizacja jako ILP: zmienne $x_{ij}, y_i$, funkcja celu,
  ograniczenia~\eqref{eq:start}--\eqref{eq:bin}, ograniczenia
  typowe~\eqref{eq:types}
& Adrian 50\%, Mikołaj 30\%, \\[-2pt]
& & Bartosz 20\% \\[4pt]

Algorytm ABC --- rdzeń metaheurystyki
& \texttt{abc\_algorithm.py}: faza employed / onlooker / scout, ruletka
  selekcji proporcjonalnej, licznik prób, reset rozwiązań,
  pętla główna i~kryterium stopu
& \textbf{Bartosz 70\%}, \\[-2pt]
& & Adrian 15\%, Mikołaj 15\% \\[4pt]

Operatory zaburzeń
& Operatory mutacji rozwiązania (add / remove / swap / 2-opt / shift)
  wraz z~rozkładem prawdopodobieństw oraz mechanizm utrzymywania
  dopuszczalności po~mutacji
& \textbf{Bartosz 65\%}, \\[-2pt]
& & Mikołaj 20\%, Adrian 15\% \\[4pt]

A\textsuperscript{*} i konstrukcja ścieżek
& Implementacja A\textsuperscript{*} na siatce 8-kierunkowej, twardy
  zakaz powtarzania krawędzi, dynamiczne blokowanie atrakcji
  spoza listy waypointów, sklejanie segmentów trasy
& Adrian 55\%, Bartosz 30\%, \\[-2pt]
& & Mikołaj 15\% \\[4pt]

Funkcja ewaluacji
& \texttt{evaluate\_path}: wartość trasy, koszt ruchu, koszt atrakcji,
  czas; leksykograficzne kryterium porównawcze rozwiązań
& Adrian 60\%, Bartosz 25\%, \\[-2pt]
& & Mikołaj 15\% \\[4pt]

Generator mapy
& \texttt{map\_generator.py}, \texttt{src/map\_generator.py}:
  procedualne generowanie wag terenu, rozmieszczanie atrakcji,
  przypisanie typów, parametry rozkładu
& \textbf{Mikołaj 65\%}, \\[-2pt]
& & Adrian 20\%, Bartosz 15\% \\[4pt]

Populacja początkowa
& \texttt{initial\_population.py}: heurystyczna konstrukcja rozwiązań
  dopuszczalnych, dobór punktów pośrednich z~uwzględnieniem
  ograniczeń typowych, kontrola czasu/budżetu
& Mikołaj 45\%, Bartosz 40\%, \\[-2pt]
& & Adrian 15\% \\[4pt]

Walidator ograniczeń
& \texttt{examples/validate\_population.py}: sprawdzanie
  spełnienia~\eqref{eq:start}--\eqref{eq:bin}, weryfikacja braku
  powtórzeń krawędzi i~ograniczeń typowych
& Adrian 65\%, Mikołaj 25\%, \\[-2pt]
& & Bartosz 10\% \\[4pt]

Parser i I/O
& \texttt{parser.py}, \texttt{load\_json.py}: konwersja JSON $\leftrightarrow$
  struktury wewnętrzne, walidacja schematu wejścia
& Mikołaj 55\%, Adrian 35\%, \\[-2pt]
& & Bartosz 10\% \\[4pt]

Wizualizacja tras
& \texttt{visualize\_paths.py}: rysowanie tras, atrakcji, punktów
  start/end, kolorowanie typów, wykresy diagnostyczne
& Mikołaj 70\%, Adrian 20\%, \\[-2pt]
& & Bartosz 10\% \\[4pt]

Aplikacja web (Streamlit)
& \texttt{app.py}, \texttt{src/path\_solver.py}, \texttt{src/ui\_io.py},
  \texttt{src/ui\_viz.py}: edytor mapy, generator, walidacja,
  tryb wyników
& Mikołaj 50\%, Adrian 35\%, \\[-2pt]
& & Bartosz 15\% \\[4pt]

Scenariusze testowe
& \texttt{scenarios/small.json}, \texttt{scenarios/medium.json}:
  parametry, dobór trudności, dane wejściowe do eksperymentów
& Adrian 45\%, Mikołaj 35\%, \\[-2pt]
& & Bartosz 20\% \\[4pt]

Dokumentacja LaTeX
& Niniejszy dokument: opis problemu, model matematyczny, opis
  algorytmu, dokumentacja użytkownika, literatura
& Adrian 40\%, Mikołaj 35\%, \\[-2pt]
& & Bartosz 25\% \\

\bottomrule
\end{longtable}

\subsection*{Sumaryczny rozkład wkładu}

Po~zsumowaniu wkładów z~poszczególnych etapów (z~wagami proporcjonalnymi
do nakładu pracy w~danym module) orientacyjny rozkład wkładu autorów
w~całość projektu przedstawia się następująco:

\begin{center}
\begin{tabular}{l c}
\toprule
\textbf{Autor} & \textbf{Łączny wkład} \\
\midrule
Bartosz Tokarz   & $\approx$ 34\% \\
Adrian Michalski & $\approx$ 34\% \\
Mikołaj Kaleta   & $\approx$ 32\% \\
\bottomrule
\end{tabular}
\end{center}

\noindent
\textbf{Komentarz do podziału.}
Bartosz Tokarz odpowiada za~rdzeń metaheurystyki ABC --- najbardziej
wymagający koncepcyjnie i~obliczeniowo komponent projektu (główna
pętla, operatory zaburzeń, mechanizm scout). Adrian Michalski skupia
się na~komponentach algorytmicznych powiązanych ze~ścieżkami
i~poprawnością rozwiązań: A\textsuperscript{*} z~ograniczeniami,
funkcja ewaluacji, walidator i~formalizacja modelu matematycznego.
Mikołaj Kaleta odpowiada za~algorytmikę po~stronie środowiska:
procedualny generator mapy, konstrukcja populacji początkowej oraz
warstwa wizualizacji i~I/O. Taki podział utrzymuje porównywalny wkład
algorytmiczny wszystkich członków zespołu, przy jednoznacznym
przypisaniu odpowiedzialności za~główny algorytm projektu.


\end{document}
